\documentclass[10pt]{article}
\usepackage[margin=0.72in]{geometry}
\usepackage{booktabs}
\usepackage{array}
\usepackage{url}
\usepackage{enumitem}
\usepackage{longtable}

\title{OptSem-C Supplemental Material}
\author{Haoyi Zhang\\
Xi'an Jiaotong-Liverpool University, Suzhou, China\\
\texttt{hyeliozhang@gmail.com}
\and
Huaijin Ran\\
Nanyang Technological University, Singapore\\
\texttt{huaijin003@e.ntu.edu.sg}}
\date{}
\begin{document}
\maketitle

\section*{Research objects}
This supplement describes the reproducibility material for \emph{OptSem-C: Benchmarking Projection Precision in Public Query-Optimizer Contracts}.  The artifact contains the grounded public-contract relation, generated query-probe benchmark, executable SQL validation corpus, external optimizer-motif representatives, projection and repair certificates, provenance records, and integrity checks.  No private datasets, account credentials, or hidden engine traces are required.

\begin{center}
\small
\begin{tabular}{lrp{0.50\linewidth}}
\toprule
Object & Count & Meaning \\
\midrule
Admitted source-linked public-contract rules & 287 & Guarded modal rules over optimizer behavior, each linked to public evidence. \\
Source-linked evidence spans & 287 & Public URL, line range, source metadata, and segment hash for every grounded rule. \\
Public source records & 26 & Documentation and system-description records used by the grounded mainline. \\
Generated SQL probes & 4,216 & Executable query probes generated from the optimizer-feature domain and validity constraints. \\
Valid feature interactions & 7,112 & Renderable optimizer-feature interactions covered by the probe generator. \\
External optimizer motifs & 90/90 & External workload and optimizer motifs covered by 71 distinct representative probes and executed on deterministic catalogs. \\
Projection-resolution subsets & 256 & Exhaustive retained-field vocabularies evaluated over the grounded comparison denominator. \\
Headline false witnesses & 498 & False equivalence witnesses across keyword, yes/no, and operator-only projections. \\
Side-balanced witnesses & 498 & False witnesses whose left and right signatures are both supported by public source records. \\
\bottomrule
\end{tabular}
\end{center}

\section*{Reproducibility paths}
\paragraph{Fast verification.}
The first command checks package coherence, frozen certificates, paper alignment, source-witness support, and file hygiene.
\begin{verbatim}
cd artifact
PYTHONDONTWRITEBYTECODE=1 ./run_mainline_checks.sh
\end{verbatim}
The expected result is that all fast checks pass and no package-integrity row fails.

\paragraph{Deep replay.}
The deep replay regenerates derived results from the grounded corpus and benchmark specifications, executes SQL validations and external motif representatives, checks projection and repair certificates, validates PDF/reference/table alignment, and rebuilds manifests.
\begin{verbatim}
cd artifact
PYTHONDONTWRITEBYTECODE=1 ./run_deep_checks.sh
\end{verbatim}
A stricter no-cache replay removes regenerable outputs before replaying the pipeline:
\begin{verbatim}
cd artifact
PYTHONDONTWRITEBYTECODE=1 ./run_from_scratch_no_cache.sh
\end{verbatim}

\section*{Certificate families}
\begin{center}
\small
\begin{tabular}{lp{0.66\linewidth}}
\toprule
Family & What is checked \\
\midrule
Grounded integrity & Rule counts, provenance coverage, source hashes, state joins, contract maps, and absence of merge conflicts. \\
Projection kernels & Exact/projection signatures, false-equivalence witnesses, information profiles, mutation projections, and exhaustive retained-field lattices. \\
Frontier antichains & Minimal safe and maximal unsafe retained-field vocabularies; monotonicity and coverage of the finite repair frontier. \\
Repair certificates & Minimum separator fields, hitting-set certificates, held-out feature stability, and layer+placement robustness. \\
Benchmark probes & Feature-domain validity, interaction coverage, SQL rendering, SQL execution, and multi-catalog SQL validation. \\
External motifs & Coverage and execution of optimizer/workload motifs drawn from recognized benchmark and query-processing families. \\
Paper alignment & Numeric claims, table sources, references, PDF integrity, manuscript style, and package consistency. \\
Evidence balance & Source breadth and left/right source support for each headline false-equivalence witness. \\
Repository quality & Architecture boundaries, package installability, cleanliness, manifest coherence, and unit tests. \\
\bottomrule
\end{tabular}
\end{center}

\section*{Main reproducibility targets}
The key numbers are: keyword projection declares 2,284 equivalences and 254 false equivalences; operator-only projection declares 2,268 equivalences and 238 false equivalences; yes/no projection declares 2,036 equivalences and 6 false equivalences; a strict exact-contract negative control declares 2,030 equivalences and zero false equivalences.  The exhaustive field-resolution lattice evaluates all 256 retained-field subsets; 44 are unsafe, and after removing action-variant identifiers the minimum safe semantic vocabularies have size two.  The layer+placement basis resolves all 498 headline witnesses.  SQL validation executes 12,648 probe-catalog pairs with zero failures across three deterministic catalog densities.

\section*{Interpretation boundary}
OptSem-C evaluates public optimizer contracts rather than hidden implementation behavior.  The corpus therefore proves what public evidence commits an engine to, not what an unobservable optimizer happens to do on a private build.  The SQL validation checks that the generated probe corpus is executable and feature-realizing; it is not a latency benchmark.  External workload motifs are used as optimizer-feature denominators rather than copied as third-party benchmark scores.  These distinctions keep the artifact aligned with the paper's central claim: coarse optimizer comparison fabricates equivalence when it drops query-processing fields, and the finite repair frontier identifies the missing semantics.
\end{document}
